\documentclass{article}
\usepackage{amsthm,amsmath}

\newtheorem{theorem}{Theorem}

\begin{document}
\begin{theorem}
Is it possible for a group to have two distinct identity elements?
\end{theorem}

\begin{proof}
Let $G$ be a group with the operation $*$. Suppose $e_1$ and $e_2$ are two
identity elements for $G$ with $e_1 \neq e_2$.

Since $e_1$ is an identity element, we have
\begin{gather*}
x*e_1=x \quad \text{and} \quad e_1*x=x
\end{gather*}
for all $x \in G$. Since $e_2 \in G$, we have $e_1*e_2=e_1$.

Similarly, since $e_2$ is an identity element, we have
$e_1*e_2=e_2$.

Thus, $e_1=e_2$. Hence, every group has a single identity element.
\end{proof}

\begin{theorem}
If $G$ is a group, and $a,b,c \in G$, then
\begin{align*}
\tag{i} ab=ac \quad &\text{implies} \quad b=c \quad \text{and}\\
\tag{ii} ba=ca \quad &\text{implies} \quad b=c
\end{align*}
\end{theorem}

\begin{proof}
Let $G$ be a group, and $a,b,c \in G$.

Now, suppose

\begin{equation}
\label{prob2:main}
ab=ac
\end{equation}

By left-multiplying \eqref{prob2:main} by $a^{-1}$, we obtain
$a^{-1}(ab)=a^{-1}(ac)$

Since $G$ is a group, by associative property, we have
\begin{align*}
(a^{-1}a)b&=(a^{-1}a)c\\
eb&=ec\\
b&=c
\end{align*}

The other case, namely $ba=ca$, follows analogously by
right-multiplying both sides by $a^{-1}$.
\end{proof}

\begin{theorem}
If $G$ is a group, and $a,b \in G$, then
\begin{align*}
\tag{i} (ab)^{-1}&=b^{-1}a^{-1} \quad \text{and}\\
\tag{ii} (a^{-1})^{-1}&=a
\end{align*}

\begin{proof}
Let $G$ be a group, and let $a,b \in G$.

(i) Since $ab \in G$, we have
\begin{align*}
(ab)(ab)^{-1}&=e\\
a^{-1}(ab)(ab)^{-1}&=a^{-1}e\\
(a^{-1}a)(b)(ab)^{-1}&=a^{-1}\\
b(ab)^{-1}&=a^{-1}\\
b^{-1}(b(ab)^{-1})&=b^{-1}a^{-1}\\
(b^{-1}b)(ab)^{-1}&=b^{-1}a^{-1}\\
(ab)^{-1}&=b^{-1}a^{-1}
\end{align*}

(ii) Since $(a^{-1})^{-1} \in G$, we have

\begin{align*}
a^{-1}(a^{-1})^{-1}&=e\\
a(a^{-1}(a^{-1})^{-1}&)=ae\\
(aa^{-1})(a^{-1})^{-1}&=a\\
(a^{-1})^{-1}&=a
\end{align*}
\end{proof}
\end{theorem}
\end{document}
